\documentclass[12pt]{article}
\usepackage[utf8]{inputenc}
\usepackage[T1]{fontenc}
\usepackage[margin=1in]{geometry}
\usepackage{setspace}
\usepackage{xcolor}
\usepackage[colorlinks=true, urlcolor=blue!60!black, linkcolor=blue!60!black]{hyperref}
\setlength{\parindent}{0em}
\setlength{\parskip}{0.45em}
\renewcommand{\baselinestretch}{1.0}
\begin{document}
\begin{center}
{\Large \textit{Griswold v. Connecticut}, 381 U.S. 479 (1965)\par}
\vspace{0.3em}
{\normalsize Case Brief --- One-Page Edition\par}
\end{center}

\textbf{Memory Jogger.} The contraception case: marital privacy lives in the ``penumbras'' of the Bill of Rights (\href{https://supreme.justia.com/cases/federal/us/381/479/}{381 U.S. 479}).

\textbf{Facts.} Connecticut made it a crime to use contraceptives. Estelle Griswold of Planned Parenthood and Dr. C. Lee Buxton opened a New Haven clinic advising married couples and were charged as accessories.

\textbf{Procedural History.} They were convicted and fined \$100 each; the Connecticut appellate courts affirmed.

\textbf{Votes.} Reversed, 7--2. Douglas wrote for the Court; Goldberg concurred; Harlan and White concurred in the judgment; Black and Stewart dissented.

\textbf{Holding.} A ban on married couples' use of contraceptives violates a fundamental right of marital privacy protected by the penumbras of specific Bill of Rights guarantees; the accessories had standing to raise it.

\textbf{Reasoning.} Disclaiming \textit{Lochner}, Douglas found ``zones of privacy'' in the First, Third, Fourth, Fifth, and Ninth Amendments, and a ban on use, rather than sale, had a ``maximum destructive impact'' on marriage. Goldberg relied on the Ninth Amendment; Harlan and White relied on due process directly. Black and Stewart found no specific provision violated, Stewart calling it ``an uncommonly silly law.''

\textbf{Examples.} \emph{Same-Side:} a ban on distribution to unmarried adults. \emph{Opposite-Side:} a prescription requirement for hormonal contraceptives. \emph{Fence-Sitter:} parental notification for minors under 16.

\textbf{Critique.} Robert Bork called the right of privacy an unprincipled leap; feminists note the marital frame excluded unmarried women.

\textbf{Key Quotation.} ``The foregoing cases suggest that specific guarantees in the Bill of Rights have penumbras, formed by emanations from those guarantees that help give them life and substance.'' --- Justice Douglas, Opinion of the Court, 381 U.S. at 484.

\end{document}
